\chapter{Introducción}\label{cap:introduccion}

AÑADIR UNA INTRODUCCIÓN

\section{Estructura de la memoria}

La memoria va a estar divididad en los siguientes capítulos:

\begin{itemize}
    \item \textbf{Capítulo 1: Introducción}
    \newline
    En este capítulo se va a dar una breve introducción del producto que se va a desarrollar en el presente Trabajo de
    Fin de Grado. Además, se van a describir los distintos capítulos en los que se dividirá la memoria.

    \item \textbf{Capítulo 2: Estado del arte}
    \newline
    En este capítulo se realizará un análisis del contexto técnico y social del problema, así como un análisis de productos
    software relacionados, herramientas existentes para su desarrollo e implementación así como librerías.

    \item \textbf{Capítulo 3: Objetivos y Alcance}
    \newline
    Aunque los objetivos hayan sido mencionados anteriormente, aquí se entrará en más detalle describiendo tanto el
    objetivo general como los objetivos específicos del trabajo.

    \item \textbf{Capítulo 4: Metdología}
    \newline
    En este capítulo se explica cómo se ha trabajado durante el desarrollo del proyecto, abarcando tanto las metodologías
    de desarrollo como posibles intervenciones con personas usuarias (entrevistas, tests de usabilidad, etc).

    \item \textbf{Capítulo 5: Análisis y requisitos}
    

    \item \textbf{Capítulo 6: Diseño del sistema y arquitectura}
    \item \textbf{Capítulo 7: Implementación}
    \item \textbf{Capítulo 8: Evaluación y experimentación}
    \item \textbf{Capítulo 9: Conclusiones y trabajos futuros}
    \item \textbf{Capítulo 10: Referencias}
    \item \textbf{Capítulo 11: Glosario}
    \newline
    Este anexo va a recoger todos aquellos términos técnicos y de relevancia que se hayan usado en el presente documento,
    con la finalidad de aclarar aquellos conceptos que puedan ser confusos o no conocidos.
\end{itemize}